\section{Introduction and main result}

The Weil explicit formula balances four sectors: an archimedean
distribution, two residual pole terms, the full prime-power measure, and a
spectral sum over the nontrivial zeros.  We study a particularly rigid test:
the interaction of two reflected translates of one fixed packet.  The
purpose is arithmetic.  No spectral ground-state or variational assertion is
made.

Let \(h\in C_c^\infty(\R)\) be the packet constructed in
Section~\ref{sec:packet}.  For \(y>1/2\), set
\[
h_y^{\mathrm R}(x)=h(x-y),\qquad
h_y^{\mathrm L}(x)=h(x+y),
\]
and
\[
h_y^\pm=\frac{h_y^{\mathrm R}\pm h_y^{\mathrm L}}{\sqrt2}.
\]
The supports of \(h_y^{\mathrm R}\) and \(h_y^{\mathrm L}\) are disjoint, so
both \(h_y^\pm\) are normalized.  The energy splitting is
\[
\Delta_h(y)=\qform[h_y^+]-\qform[h_y^-].
\]

Write
\[
K_h=h*h,\qquad
\Lh(w)=\int_{-1}^{1}K_h(r)e^{wr}\,dr
\]
and define the all-prime-power error measure
\[
dE(x)=\sum_{n\ge2}\Lambda(n)\delta_n-dx.
\]
The centered arithmetic interaction is
\[
\Ih(y)=
-\int_0^\infty x^{-1/2}K_h(2y-\log x)\,dE(x).
\]

\begin{theorem}[Exact expansion and unconditional oscillation]
\label{thm:main}
For every \(y>1/2\),
\begin{align}
\Ih(y)
={}&
\sum_{\rho}m_\rho
\Lh\left(\rho-\frac12\right)e^{2(\rho-1/2)y}
\notag\\
&+
\sum_{n=1}^{\infty}
\Lh\left(2n+\frac12\right)e^{-(4n+1)y},
\label{eq:I-expansion}
\end{align}
where \(\rho\) ranges over the distinct nontrivial zeros of \(\zeta\) and
\(m_\rho\) is the multiplicity.  Both series converge absolutely, and the
nontrivial-zero series may be differentiated term by term any fixed number
of times locally uniformly in \(y\).

The archimedean and residual pole terms cancel the second series exactly.
Consequently
\begin{equation}
\qform(h_y^{\mathrm R},h_y^{\mathrm L})
=
\sum_\rho m_\rho
\Lh\left(\rho-\frac12\right)e^{2(\rho-1/2)y},
\label{eq:zero-only}
\end{equation}
and
\begin{equation}
\Delta_h(y)=
2\sum_\rho m_\rho
\Lh\left(\rho-\frac12\right)e^{2(\rho-1/2)y}.
\label{eq:delta-zero-only}
\end{equation}
The coefficient at \(\rho\) is
\[
A_h(\rho)=m_\rho\Lh\left(\rho-\frac12\right),
\]
and it vanishes exactly if
\[
\rho=\frac12+4\pi i k,\qquad k\in\Z\setminus\{0\}.
\]

Unconditionally, for every \(Y>0\), there exist
\(y_+,y_-,u_+,u_->Y\) such that
\[
\Ih(y_+)>0,\qquad \Ih(y_-)<0,
\]
and
\[
\Delta_h(u_+)>0,\qquad \Delta_h(u_-)<0.
\]
Thus each function has infinitely many sign-crossing intervals.
\end{theorem}

The sign convention has a direct but limited interpretation:
\(\Delta_h(y)<0\) means that the even packet has lower energy than the odd
packet at that separation, while \(\Delta_h(y)>0\) means the reverse.
Theorem~\ref{thm:main} therefore proves that neither parity is eventually
preferred for this one packet.

The proof has four independent components.  First, exact polarization fixes
the factor \(2\).  Second, the pole and prime-power main terms are combined
before estimation.  Third, a contour displacement produces
\eqref{eq:I-expansion}, including every trivial zero.  Finally, an exact
packet-transform zero set and Conrey's theorem produce visible nonreal poles;
Landau's theorem then rules out either eventual sign.
